\section*{Introduction}

Genome-wide association studies have identified tens of thousands of
trait-associated loci\cite{visscher2017years,buniello2019nhgri},
but resolving the causal variants that drive each signal remains the
rate-limiting step between statistical discovery and mechanistic
biology. Because nearby variants co-segregate through linkage
disequilibrium (LD), a single causal variant is typically
indistinguishable from hundreds of non-causal proxies on statistical
evidence alone. Fine-mapping procedures aim to break this
degeneracy\cite{schaid2018genome,ulirsch2021fine}; their accuracy
directly constrains downstream experimental follow-up and the
trans-ancestry portability of polygenic scores.

The challenge has different shapes across the regimes in which
fine-mapping is most needed. In human medical genetics, biobank-scale
sample sizes deliver high signal-to-noise but LD is dense and
ancestry-imbalanced, so the same lead variant can map to tightly
correlated proxy clusters of $10^2$--$10^3$ variants in non-European
panels even when the European credible set resolves to a singleton.
In the model organisms used to interrogate trait architecture under
controlled conditions --- inbred yeast and \emph{Arabidopsis} 1001G ---
LD is exceptionally long, so single-locus credible sets routinely
exceed $10^3$ variants and statistical resolution alone cannot
isolate causal variants. In agricultural germplasm such as the
3{,}000 Rice Genomes panel\cite{wang2018genomic}, accessions span
between-population $F_{ST}\!\approx\!0.5$, most traits have only
ordinal phenotypic scoring, and curated causal-gene catalogues exist
for only a few well-studied phenotypes\cite{ren2023rice}. A method
that improves fine-mapping must therefore travel between these
regimes without re-engineering for each.

Modern fine-mapping is dominated by Bayesian variable-selection
methods operating on a per-locus genotype matrix or summary-statistics
vector---SuSiE\cite{wang2020simple}, FINEMAP\cite{benner2016finemap},
PolyFun\cite{weissbrod2020polyfun},
PAINTOR\cite{kichaev2014integrating}---with two recent extensions
that add an infinitesimal random-effect background
(SuSiE-inf and FINEMAP-inf)\cite{cui2024improving} and a genome-wide
multi-component Bayesian mixture
(SBayesRC)\cite{zheng2024leveraging} that
Wu~\emph{et~al.}~2026\cite{wu2026genome} used as the substrate for a
genome-wide fine-mapping framework. All share a common design choice:
biological annotations, where used at all, enter as \emph{flat
per-variant scalars}---averaged across tissues, collapsed across
pathways, and stripped of their relational structure. A variant's
connection to a particular gene in a particular tissue, participating
in a particular pathway coupled via protein--protein interactions to
further functional neighbourhoods, is flattened to a single weight.
Constructing that scalar is itself a curation step that scales poorly
to non-model species --- including crops, livestock and model
organisms --- where functional annotations are sparse, fragmented
and lack a pre-harmonised category catalogue.

This flattening is the limiting approximation. Multi-omics resources
released since 2020---the Genotype-Tissue Expression (GTEx) project
v8 with 43~million tissue-specific expression quantitative trait
loci (eQTLs)\cite{gtex2020}, STRING~v12 with 230{,}000 high-confidence
protein--protein interactions (PPIs) at combined
score~$\geq$~700\cite{szklarczyk2023string}, and the ENCODE Project's
candidate cis-regulatory elements (cCREs) with 370{,}000
entries\cite{encode2020expanded}---carry exactly the relational
biology that flat priors discard. When signal is strong, statistical evidence alone resolves
the causal variant and relational context is optional; but this is
not the regime that matters most for novel discovery. Sub-genome-wide-
significant loci, low-frequency effects and non-European ancestries
all sit in the weak-signal regime, where annotation priors dominate
the posterior. The question we address is whether fine-mapping
methods that represent multi-omics context \emph{as a graph, not as a
vector} can decisively outperform flat-prior methods in that regime.

Answering that question requires a computational substrate in which
multi-omics structure is \emph{primary data}, not reconstructed on
demand from files. Graph databases provide exactly
this\cite{robinson2015graph}: typed nodes (variant, gene, pathway,
tissue, regulatory element) and labelled edges (consequence, eQTL,
pathway membership, protein--protein interaction) are first-class
storage objects; index-free adjacency makes traversing a
relationship an $O(1)$ operation regardless of database size; and
new annotation layers register as additional edge types without
restructuring the underlying schema. These properties are absent
from the file-based pipelines assembled from variant call format
(VCF) files, browser extensible data (BED) tracks, general feature
format (GTF) annotations, and pathway tables glued together by
\texttt{bedtools}/\texttt{tabix} lookups. Encoding the biology
directly as a graph reframes fine-mapping as a Bayesian computation
\emph{on the graph itself}---not as a regression that happens to
have annotation weights attached---and collapses what would
otherwise be a multi-file data-integration task into a single
database traversal that can be iterated thousands of times per
second.

Doing so changes the fine-mapping
accuracy--speed frontier. GraphGWAS is a fine-mapping platform built
on a graph database\cite{robinson2015graph} into which users import
their own genotype, summary-statistics and multi-omics data; the
schema and algorithms are database-agnostic and the present
reference implementation runs on Neo4j~5.26 Community Edition. The
demonstration loadout used in this paper covers 70.7~M variants from
the 1000 Genomes Phase~3 high-coverage release, 20{,}092 GENCODE~v47
genes\cite{frankish2023gencode}, 43.2~M GTEx~v8 eQTL edges,
230{,}850 STRING interactions and 370{,}000 ENCODE regulatory
elements. Its central methods ---
hierarchical belief propagation (HBP), with proved geometric
convergence via a Banach contraction argument, and graph-augmented
fine-mapping (GAFM), with a proof that the causal variant ranks
first under mild LD-decay assumptions --- carry the full
variant~$\rightarrow$~gene$\rightarrow$~pathway factor graph with
PPI coupling through every update. On 90 simulated chromosome~22
loci, HBP matches SuSiE and FINEMAP at 5--40$\times$ the speed; on
79 weak-signal tissue-specific eQTL simulations, GAFM beats SuSiE
27--2 head-to-head (sign-test $p < 10^{-5}$). Posterior inclusion
probabilities (PIPs) are calibrated with 0\% null false-positive
rate (FPR) across 100 simulations. A
sumstats-only entry path consumes Pan-UK Biobank\cite{karczewski2024pan}
directly, fine-mapping across four ancestries
(EUR $N{=}420{,}531$; CSA, AFR, EAS) and recovering canonical
lipid genes \emph{LDLR}, \emph{APOE}, \emph{LPL}, \emph{GCKR},
\emph{ANGPTL3} at single-variant resolution. The same code
fine-maps yeast\cite{peter2018genome},
\emph{Arabidopsis}\cite{alonso20161001} and the 3{,}000 Rice
Genomes panel without algorithmic change. Across 692 leads in
the four species, a direct intervention experiment shows that the
operative quantity is per-variant heterogeneity of the prior, not
annotation density --- adding even a coarse regulatory-element
flag to the otherwise dense uniform human cache flips HBP
narrower than flat-prior GAFM from 0\% to 88\% of leads.
